\section{A coefficient-free oracle for univariate shallow networks}
\label{sec:univariate-coefficient-free}
Let $X\in[-1,1]$ and suppose that, for every interval $I\subseteq[-1,1]$,
\begin{equation}
 \Pr(X\in I)\le\kappa |I|^\alpha,
 \qquad \kappa\ge1,\qquad 0<\alpha\le1.
 \label{eq:univariate-regularity}
\end{equation}
A density bounded by $\kappa$ gives the case $\alpha=1$. The constants
below do not depend on the original teacher coefficients.

\begin{lemma}[Bounded-slope approximation of a bounded spline]
\label{lem:bounded-slope-spline}
Let $f:[-1,1]\to[-1,1]$ be continuous and piecewise affine, with at most
$K$ interior breakpoints. For every $0<\delta\le1$ there is a continuous
piecewise-affine $q:[-1,1]\to[-1,1]$ with at most $2K$ interior
breakpoints and Lipschitz constant at most $2/\delta$ such that
\begin{equation}
 \E|q(X)-f(X)|\le6\kappa K\delta^\alpha.
 \label{eq:spline-approximation}
\end{equation}
Moreover, $q$ has a one-hidden-layer ReLU implementation, with a
constant input, at most $2K+2$ hidden neurons, and every coefficient
bounded in absolute value by $4/\delta$.
\end{lemma}
\begin{proof}
If $K=0$, take $q=f$. A bounded affine function on $[-1,1]$ has slope
at most one, and the representation below uses two neurons.
Otherwise list the actual breakpoints as
$-1<t_1<\cdots<t_k<1$, with $k\le K$. Partition them into maximal
consecutive clusters by placing two successive breakpoints in the
same cluster exactly when their gap is at most $3\delta$.
For a cluster with first and last breakpoints $a,b$, set
\[
 I_{a,b}=[\max\{-1,a-\delta\},\min\{1,b+\delta\}].
\]
The resulting intervals are disjoint, and successive such intervals
are separated by more than $\delta$. Define $q=f$ outside their
interiors, and interpolate linearly between the two endpoint values
of $f$ on each interval. This preserves continuity and the bound
$|q|\le1$. All original breakpoints lie inside these intervals, so
$q$ has at most two breakpoints per cluster.

Each replacement interval has length at least $\delta$, hence its
chord has slope at most $2/\delta$. Every complementary interval
between two replacement intervals has length greater than $\delta$;
on it, $f$ is affine and bounded, giving the same slope bound.
If a nonempty complementary interval touches the left endpoint
$-1$, then the first original breakpoint lies more than $\delta$
from $-1$. The affine segment extends all the way to that breakpoint,
so its slope is again at most $2/\delta$ even if the retained
complementary interval itself is short. The right boundary is
identical. Thus $q$ is globally $2/\delta$-Lipschitz.

A cluster containing $k_j$ breakpoints has replacement-interval
length at most
$3\delta(k_j-1)+2\delta\le3k_j\delta$.
Since $|q-f|\le2$ and the functions agree elsewhere,
\[
 \E|q(X)-f(X)|
 \le 2\kappa\sum_j(3k_j\delta)^\alpha
 \le 6\kappa\delta^\alpha\sum_j k_j
 \le 6\kappa K\delta^\alpha.
\]
Here $3^\alpha\le3$ and $k_j^\alpha\le k_j$.

For completeness, if $u_1<\cdots<u_r$ are the new interior
breakpoints and $s_0,\ldots,s_r$ the slopes, use
\[
 q(x)=q(-1)\operatorname{ReLU}(1)
       +s_0\operatorname{ReLU}(x+1)
       +\sum_{j=1}^r(s_j-s_{j-1})\operatorname{ReLU}(x-u_j).
\]
The identity holds throughout $[-1,1]$. All input coefficients and
biases are at most one, $|q(-1)|\le1$, $|s_0|\le2/\delta$, and every
slope jump is at most $4/\delta$. This proves the implementation claim.
\end{proof}

\begin{corollary}[Compactifying a clipped univariate ReLU network]
\label{cor:univariate-compiler}
Let $g$ be any one-hidden-layer univariate ReLU network of width
$C\ge1$, with arbitrary finite real coefficients and biases. Let $s$
clip to $[-1,1]$. For every $0<\zeta\le1$ there is a one-hidden-layer
ReLU network $q$ satisfying
\begin{equation}
 |q(x)|\le1\quad(x\in[-1,1]),\qquad
 \E|q(X)-s(g(X))|\le\zeta,
 \label{eq:univariate-compiler-error}
\end{equation}
whose width and coefficient cap obey
\begin{equation}
 m\le6C+6,\qquad
 B\le4\left(\frac{18\kappa(C+1)}{\zeta}\right)^{1/\alpha}.
 \label{eq:univariate-compiler-resources}
\end{equation}
It has a constant input and no internal biases, and its nonnegative
hidden features can be relayed by identity matrices through any
additional hidden layers.
\end{corollary}
\begin{proof}
The original network is affine on at most $C+1$ intervals. Clipping
can add at most two breakpoints on each such interval, so $s\circ g$
has at most $K=3C+2$ interior breakpoints. Apply
Lemma~\ref{lem:bounded-slope-spline} with
$\delta=[\zeta/(18\kappa(C+1))]^{1/\alpha}\le1$.
Then $6\kappa K\delta^\alpha\le\zeta$ and
$2K+2\le6C+6$. A ReLU identity layer preserves the nonnegative
hidden vector exactly.
\end{proof}

\begin{theorem}[A sharp continuous-domain width-only oracle]
\label{thm:univariate-levy-oracle}
Use the independent-coordinate L\'evy fan-in prior with the fixed small
intensity required by the stable-padding comparison theorem. Fix one
learner width $N$, one hidden depth $H\ge1$, ReLU hidden activation,
and clipping to $[-1,1]$ in both the likelihood and prediction.
Assume \eqref{eq:univariate-regularity},
$f_0(x)=\E[Y\mid X=x]\in[-1,1]$, and conditionally centered
$\sigma^2$-sub-Gaussian label noise. Put $s_\sigma=1+\sigma^2$.
The ordinary all-data posterior with likelihood
\[
 d\Pi_{S_n}(w)\ \propto\
 \exp\!\left[-\frac1{256s_\sigma}
          \sum_{i=1}^n(f_w(X_i)-Y_i)^2\right]d\pi_N(w)
\]
and posterior mean $\bar f_{S_n}$ satisfies, for a numerical $A$,
\begin{equation}
 \E_{S_n}R(\bar f_{S_n})
 \le
 \inf_{\substack{C\ge1:\,6C+6\le N}}
 \left\{\inf_{g\in\mathcal G_C}R(g)
 +\frac{A s_\sigma(C+1)H}{\alpha(n+1)}
       \log\!\bigl[e\kappa(C+1)(H+1)(n+1)\bigr]\right\},
 \label{eq:univariate-width-oracle}
\end{equation}
where $\mathcal G_C$ is the class of one-hidden-layer univariate ReLU
networks of width at most $C$, with no coefficient restriction. The
same fixed posterior is used for every comparator in this infimum.
If the displayed index set is empty, the theorem makes no comparator
claim. No prior over network sizes or posterior averaging over sample
sizes is used.
\end{theorem}
\begin{proof}
Fix $C$ with $6C+6\le N$ and $g\in\mathcal G_C$.
Clipping $g$ cannot increase squared risk because $f_0\in[-1,1]$.
Apply Corollary~\ref{cor:univariate-compiler} with
$\zeta=(n+1)^{-1}$ to get a bounded width-$m$ comparator $q$.
For two functions taking values in $[-1,1]$, population squared risks
differ by at most four times their $L^1(P_X)$ distance; this follows
by conditioning on $X$ and cancelling the common label-noise variance.
Consequently
\[
 R(q)\le R(g)+\frac4{n+1}.
\]
Its one-hidden-layer rooted formula has size at most $m$. Identity
padding to depth $H$ increases this to at most $mH$. The audited
stable-padding comparison theorem gives the all-data expectation
remainder
\[
 \frac{A s_\sigma(mH+1)}{n+1}
 \left[2\log\{A(m+1)(B+8)\}
       +\log\{A(mH+1)H(B+8)(n+1)\}\right].
\]
Substitute \eqref{eq:univariate-compiler-resources} and
$m\le6C+6$. Since $0<\alpha\le1$, this is bounded by the remainder
in \eqref{eq:univariate-width-oracle}, after enlarging $A$ to absorb
$4/(n+1)$. Every comparison law is deterministic given the comparator
and input distribution. The expectation oracle holds separately for
every such law and uses the same posterior, so one may take the
infimum over all $C$ and $g$ without a confidence union or a size prior.
\end{proof}

\paragraph{Interpretation.}
For a bounded-density one-dimensional input law and fixed learner
depth, the excess-risk cost is $\widetilde O((1+\sigma^2)C/n)$,
independent of the original teacher coefficients and the learner
width. Additional learner depth costs only a factor $H$ and logarithms.
The anti-concentration assumption matters: the proof permits changes
on short intervals around tightly packed breakpoints, and controls
their probability through \eqref{eq:univariate-regularity}. The theorem
is an expectation bound for the ordinary final posterior mean. It
does not claim that a single posterior draw has the same guarantee.
